\chapter{The PED Numerical System}

\section{Numerical Stability and the self}
In the PED conceptual framework, numbers were not abstract mathematical entities but distances from the agential self. The earliest counting system was a binary-quinary hybrid, rooted in the topography of the body.

\section{The Cardinal Reconstructions}
We provide the systematic derivation of the numerals 1, 2, and 10, demonstrating the exact application of the Glottalic Shift Law across the functional category of numerals.

\begin{table}[h]
\centering
\begin{tabular}{@{}lllll@{}} \toprule
\textbf{Numeral} & \textbf{PED Root} & \textbf{Vedic} & \textbf{Tamil} & \textbf{Tibetan} \\ \midrule
1 & *oK'n & oin- & on- & it \\
2 & *T'uH & du- & tu- & ni \\
10 & *T'eK' & dek- & tek- & gip \\ \bottomrule
\end{tabular}
\caption{Reconstructed Numerals (1, 2, 10).}
\end{table}

\section{Structural Proof: The 'Two' Cluster}
The alignment of the dental ejective \textit{*T'} with Vedic \textit{d} and Tamil \textit{t} in the numeral for 'two' provides a degree of systematicity that cannot be attributed to chance. The probability of this specific sound shift occurring across independent functional categories (e.g., 'tree', 'indicate', 'water', and 'two') is vanishingly small, satisfying the requirement for non-circular proof.
